\section{Related Work}

\label{sec:related}

\subsection{Theorem Search Engines and Infrastructure}

\textbf{TheoremSearch}~\cite{theoremsearch2026} indexes 9.2 million statements extracted from arXiv and seven additional sources, represents each theorem with a short natural language description for embedding, and exposes a public API without authentication. Its documented motivation (arXiv withdrawals due to duplication) is the same that motivates our experiment with withdrawn papers. TheoremSearch finds similar statements; AViD adds the verdict layer.

\textbf{Matlas}~\cite{matlas2026} extracts 8.07 million statements from 435,000 peer-reviewed articles spanning 1826 to 2025, sourced from 180 journals selected by ICM citation criteria, plus 1,900 textbooks. The two sources are complementary in temporal coverage: TheoremSearch covers arXiv since 1991, Matlas covers peer-reviewed journals since 1826. AViD consumes them as the two branches of its informal corpus, and section~\ref{sec:experimentos} reports how far that combined coverage reaches.

The Mathlib search engine space is saturated. \textbf{Leandex}~\cite{leandex2026}, Project Numina's semantic search engine, offers search over Lean~4 declarations; \textbf{Lean Finder}~\cite{leanfinder2025} incorporates user intent into semantic search over Mathlib. AViD does not compete in that space: it consumes one of them as infrastructure (Leandex is the backend for D1~C$_F$) and adds the decision layer that none offers. A recent survey on mathematical AI~\cite{survey2026} cites statement search engines as infrastructure for locating already established theorems, which situates the problem AViD addresses as recognized by the field.

Two recent works are relevant as infrastructure, not as competitors. \textbf{TheoremGraph}~\cite{theoremgraph2026} builds a statement-level dependency graph that unifies the informal (11.7M theorem-type environments from arXiv) and the formal (LeanGraph, 388,105 Lean~4 nodes); we evaluated using its graph dependencies to measure distance between proofs in D3, but ultimately stayed with the direct premise set of each proof. In a complementary direction, \textbf{COMPOSE}~\cite{compose2026} predicts plausible future statements by combining an article's citation graph with Mathlib's formal dependencies, generating conjectures rather than arbitrating their novelty.

\subsection{Novelty Filters in Conjecture Generation}

\textbf{LeanConjecturer}~\cite{leanconjecturer2025} filters novelty with \texttt{exact?} against Mathlib and non-triviality with \texttt{aesop}: it generated 12,289 conjectures from 40 seed files, of which 3,776 proved syntactically valid and non-trivial. These are exactly the mechanisms that AViD implements as D1 (\texttt{exact?} as fallback for C$_F$) and D2 (\texttt{aesop} in \texttt{T\_AUTO}). AViD's contribution is not inventing these filters but composing them into a decision tree that issues a verdict, adding D3 and an informal branch with temporal filter, and evaluating them against external ground truth. Work on pruning redundant conjectures~\cite{mining2024} and synthetic generation via forward reasoning~\cite{synthetic2024} illustrates the same need from the generation side.

\subsection{Proof Identity and Similarity}

The Jaccard distance over premise sets that AViD uses in D3 borrows its representation from the premise selection literature, where a theorem is characterized by the premises of its proof: MaSh~\cite{kuhlwein2013}, \textbf{Kaliszyk and Urban}~\cite{kaliszyk2014} with IDF-weighted $k$-NN, and \textbf{Magnushammer}~\cite{magnushammer2024} with transformers. But the task is the inverse: those systems \textit{select} premises to construct a proof from a statement, whereas D3 uses the set of premises \textit{already employed} as a fingerprint to compare finished proofs. The only directly inherited component is the filtering: the \textit{math filter} of \textbf{Piotrowski et al.}~\cite{piotrowski2023}, which uses Mathlib names as a whitelist to discard technical lemmas, is the same recipe applied by the filters prior to D3's Jaccard.

The Mathlib network analysis by Li et al.~\cite{network2026}, over a graph of 308,129 declarations and 8.4 million edges, shows that the most cited lemmas in the library are technical infrastructure rather than deep mathematics: equality reflexivity (\texttt{Eq.refl}) is the second highest in-degree node, with 69,580 citations, while the Chinese Remainder Theorem does not even appear among the top hundred. That is, centrality in the graph measures how ubiquitous a lemma is, not how much mathematical content it contributes. This justifies D3's Filter~1: before computing the Jaccard between two proofs, AViD discards premises from the \texttt{Init.} and \texttt{Lean.} namespaces, which are precisely those omnipresent technical lemmas. Without the filter, two proofs would appear similar merely for sharing the basic machinery of the language; with it, the Jaccard reflects the premises that truly distinguish mathematical content.

\textbf{Yoo}~\cite{yoo2025} represents theorems as \textit{proof vectors} over axiom systems and compares them with cosine, Euclidean, or Jaccard similarity. It is the closest work in conceptual tool but with opposite purpose: the Atlas organizes theorems by structural similarity; AViD uses that similarity to decide whether a proof is redundant. \textbf{Huch}~\cite{huch2022} documents the scale-free in-degree distribution in the AFP, relevant for future work with IDF-weighted premises.

\subsection{Correctness versus Novelty}

Existing paper-to-Lean pipelines verify that a proof is correct, not that it is new; novelty verification is delegated to human reviewers. AViD is explicitly that delegated layer. A complete automated review system would require both dimensions, and section~\ref{sec:experimentos} shows that the first is not a sufficient condition for the second: a Lean file can compile without errors and yet not be a faithful formalization of the original statement.

\subsection{Positioning}

\begin{table}[htbp]
\centering
\small
\setlength{\tabcolsep}{4pt}
\begin{tabular}{lccccc}
\toprule
System & Existence & Triviality & Proof dist. & Informal corpus & Verdict \\
\midrule
Kasaura et al.~\cite{kasaura2025}      & Mathlib            & --- & --- & --- & yes/no \\
LeanConjecturer~\cite{leanconjecturer2025} & \texttt{exact?} & \texttt{aesop} & --- & --- & --- \\
TheoremSearch~\cite{theoremsearch2026} & 9.2M statements    & --- & --- & \checkmark & --- \\
Matlas~\cite{matlas2026}               & 8.07M statements   & --- & --- & \checkmark & --- \\
Atlas~\cite{yoo2025}                   & ---                & --- & Jaccard & --- & --- \\
AViD (this work)                      & Mathlib + informal & 6 tactics & Jaccard$^{*}$ & TS + Matlas & 8 verdicts \\
\bottomrule
\end{tabular}
\caption{Positioning of AViD with respect to prior work. Prior works build infrastructure pieces; AViD composes them into a pipeline that goes from LaTeX source to verdict. $^{*}$D3 is implemented and validated in isolation (section~\ref{sec:instrumentos}), but is not exercised within the decision tree in the reported experiments.}
\label{tab:posicionamiento}
\end{table}

The difference lies not in any individual cell but in the composition: prior works build infrastructure pieces---search, formalization, similarity---; AViD assembles them into a pipeline that receives a LaTeX paper and returns a verdict.